\documentclass[11pt,a4paper]{article}

\input{../preambule}

\begin{document}

\title{Semaine 4}
\author{}
\date{12/10 - 19/10 }

\maketitle


\begin{exercise}
Soit $I$ un intervalle, $f$ continue sur $I$ à valeurs dans $[a,b]$ avec $a\lt b\in\R$.\\
Montrer que $\displaystyle\int_If^2\le-ab|I|$.
\end{exercise}

\begin{exercise}
Soit $f$ une fonction continue sur $[0,1]$.\\
Calculer $\lim\limits_{n\to+\infty}\int_0^1f(t^n)dt$
\end{exercise}

\begin{exercise}
Montrer que pour tout $t\gt0, \arctan(t)\ge\frac{t}{1+t^2}$
\end{exercise}

\begin{exercise}
Soit $z_1,\ldots, z_n\in\C$. Montrer qu'il existe $I\subset\intset{1,n}$ tel que $\abs{\sum_{i\in I}z_i}\ge\frac{1}{6}\sum_{i=1}^n\abs{z_i}$
\end{exercise}

\begin{exercise}
Soit $f\in C^0(\R^+,\R^+)$ telle qu'il existe $k\ge0$ tel que pour tout $x\gt0, f(x)\le k\int_0^xf(t)dt$.
Montrer que $f$ est nulle.
\end{exercise}

\begin{exercise}
Soit $f$ une fonction continue sur $[a,b]$ et $n\in\N$ telle que $\forall k\in\intset{0,n},\int_a^bt^kf(t)dt=0$.
Montrer que $f$ s'annule au moins $n+1$ fois sur $[a,b]$.
\end{exercise}





\end{document}